\exercisesection{}

除另有说明外，在本节的各习题中，E 表示一个复希尔伯特空间。

\begin{enumerate}
\def\labelenumi{\arabic{enumi})}
\tightlist
\item
  设 \(u\) 是 E 的自同态，其像在 E 中是闭的。设 \(q\) 是 E 的正交投影算子，并设 \(v\) 是 E 的像的自同态 \(x \mapsto q(u(x))\) 。
\end{enumerate}

\begin{enumerate}
\def\labelenumi{\alph{enumi})}
\item
  \(v\) 的伴随是 E 的像的自同态 \(x \mapsto q(u^{*}(x))\) 。
\item
  若 \(u\) 是埃尔米特的，则 \(v\) 是埃尔米特的。
\item
  有可能 \(u\) 是正规的而 \(v\) 不是正规的。
\end{enumerate}

\begin{enumerate}
\def\labelenumi{\arabic{enumi})}
\setcounter{enumi}{1}
\item
  直接利用谱定理（IV, p. 193 的定理 1 的推论）证明 IV, p. 149 的定理 1。
\item
  设 \(S \subset \mathbf{C}\) 是 \(\mathbf{C}\) 的一个紧子集。设 \(\varphi \colon \mathcal{C}(S) \to \mathcal{L}(E)\) 是对合代数的一个连续的含单位元态射。存在唯一的正规自同态 \(u\) ，它作用在 \(E\) 上，谱为 \(S\)，并且使得 \(\varphi(f) = f(u)\) 对所有 \(f \in \mathcal{L}^{\infty}(S)\) 成立。
\item
  设 \(u \in \mathcal{L}(\mathrm{E})\) 是正规的。\(u\) 的谱的凸包是 \(\mathbf{C}\) 中数值值域 \(\iota(u)\) （相应于 \(u\) ；参见 I, p. 135 的定义 2）的闭包。
\item
  设 \(u \in \mathcal{L}(\mathrm{E})\) 满足 \(\| u \| \leqslant 1\) 。对每个 \(\lambda \in \mathbf{C}\) ，用 \(p_{\lambda}\) 表示 E 到 \(u\) 的对应于 \(\lambda\) 的特征空间上的正交投影算子。对 E 中所有的 \(x\) 与 \(y\)，我们有
\end{enumerate}

\[
\lim_{N\to +\infty}\frac{1}{N}\sum_{n = 1}^{N}|\langle x \mid u^{n}(y)\rangle |^{2} = \sum_{\substack{\lambda \in \mathbf{C}\\ |\lambda | = 1}}|\langle p_{\lambda}(y) \mid p_{\lambda}(x)\rangle^{2}|\\ = \sum_{\substack{\lambda \in \mathbf{C}\\ |\lambda | = 1}}|\langle y \mid p_{\lambda}(x)\rangle |^{2}
\]

（可以利用 II, p. 300 的习题 64）。

\begin{enumerate}
\def\labelenumi{\arabic{enumi})}
\setcounter{enumi}{5}
\tightlist
\item
  设 E 是预希尔伯特空间，F 是巴拿赫空间，\(u: E \rightarrow F\) 是线性映射。
\end{enumerate}

\begin{enumerate}
\def\labelenumi{\alph{enumi})}
\item
  证明：u 是紧线性映射当且仅当 \(\|u(e_{n})\|\to0\) 对 E 的每个标准正交序列 \((e_{n})\) 成立。（为证明该条件是充分的，利用 TG, IX, p. 20 的命题 14。）
\item
  假设 E 是可分希尔伯特空间且 F = E。为使存在 E 的标准正交基 \((e_{n})\) 使得 \(\|u(e_{n})\| \to 0\)，必须且只需 u 不是弗雷德霍姆算子。
\item
  假设 E 是希尔伯特空间，并且 \(\langle u(e_{n}), e_{n} \rangle \to 0\) 对 H 的每个标准正交序列 \((e_{n})\) 成立。证明 u 是紧线性映射。
\end{enumerate}

\begin{enumerate}
\def\labelenumi{\arabic{enumi})}
\setcounter{enumi}{6}
\tightlist
\item
  设 \(u\) 与 \(v\) 是 E 的连续自同态。
\end{enumerate}

\begin{enumerate}
\def\labelenumi{\alph{enumi})}
\item
  证明 \(\operatorname{Im}(u) + \operatorname{Im}(v) = \operatorname{Im}((u \circ u^* + v \circ v^*)^{1/2})\) 。
\item
  若 \(u\) 是正的且具有闭像，证明 \(\operatorname{Im}(u) = \operatorname{Im}(u^{1/2})\) 。
\item
  若 \(u, v\) 是正的且具有闭像，并且 \(u + v\) 也具有闭像，则我们有 \(\operatorname{Im}(u + v) = \operatorname{Im}(u) + \operatorname{Im}(v)\) 。
\end{enumerate}

\begin{enumerate}
\def\labelenumi{\arabic{enumi})}
\setcounter{enumi}{7}
\tightlist
\item
  设 F 是复希尔伯特空间，\(u \colon \mathrm{E} \to \mathrm{F}\) 是具有闭像的连续线性映射。
\end{enumerate}

\begin{enumerate}
\def\labelenumi{\alph{enumi})}
\tightlist
\item
  证明存在唯一的连续线性映射 \(v \in \mathcal{L}(\mathrm{F};\mathrm{E})\) 满足下列条件：
\end{enumerate}

\begin{enumerate}
\def\labelenumi{(\roman{enumi})}
\item
  \(u \circ v \circ u = u\) 且 \(v \circ u \circ v = v\) ；
\item
  \(u \circ v\) 与 \(v \circ u\) 是自伴的。
\end{enumerate}

我们把这一映射记作 \(u^{\dagger}\) （穆尔–彭罗斯伪逆）。

\begin{enumerate}
\def\labelenumi{\alph{enumi})}
\setcounter{enumi}{1}
\tightlist
\item
  证明 \(\operatorname{Im}(u) = \operatorname{Im}(u \circ u^{\dagger})\)。证明：对 \(x \in E\) 与 \(y \in \operatorname{Im}(u)\) ，条件 \(u(x) = y\) 与 \(x - u^{\dagger}(y) \in \operatorname{Im}(1_{\operatorname{E}} - u^{\dagger} \circ u)\) 等价。
\end{enumerate}

\begin{enumerate}
\def\labelenumi{\arabic{enumi})}
\setcounter{enumi}{8}
\tightlist
\item
  设 \(u\) 与 \(v\) 是 E 的具有闭像的自同态。若 \(u + v\) 也具有闭像，我们令 \(u: v = u \circ (u + v)^{\dagger} \circ v\) ，其中 \((u + v)^{\dagger}\) 是 \(u + v\) 的穆尔–彭罗斯伪逆。
\end{enumerate}

\begin{enumerate}
\def\labelenumi{\alph{enumi})}
\item
  证明 \(u:v\) 是自伴的。
\item
  我们有 \(u:v=v:u\) 且 \(\operatorname{Im}(u:v)=\operatorname{Im}(u)\cap\operatorname{Im}(v)\) 。
\item
  假设 u 与 v 是正交投影算子。证明 2u : v 是以 \(\operatorname{Im}(u) \cap \operatorname{Im}(v)\) 为像的正交投影算子。
\end{enumerate}

\begin{enumerate}
\def\labelenumi{\arabic{enumi})}
\setcounter{enumi}{9}
\tightlist
\item
  E 的连续自同态 \(u\) 称为亚正规的，如果 \(u^{*} \circ u \geqslant u \circ u^{*}\)；称为仿正规的，如果 \(\|u^{2}(x)\|\|x\| \geqslant \|u(x)\|^{2}\) 对所有 \(x \in E\) 成立。
\end{enumerate}

\begin{enumerate}
\def\labelenumi{\alph{enumi})}
\item
  若 \(u\) 是仿正规的，则 \(u^n\) 对每个 \(n \in \mathbf{N}\) 都是仿正规的。
\item
  若 u 是仿正规的且可逆，则 \(u^{-1}\) 是仿正规的。
\item
  亚正规自同态是仿正规的。
\item
  设 u 是谱包含在单位圆周中的仿正规自同态；则 u 是酉的。
\item
  若 \(u\) 是仿正规自同态，则它的范数 \(\| u\|\) 等于它的谱半径 \(\varrho(u)\) 。
\item
  我们假设 u 是仿正规的。
\end{enumerate}

\begin{enumerate}
\def\labelenumi{\arabic{enumi})}
\setcounter{enumi}{10}
\tightlist
\item
  设 u 是 E 的连续自同态。
\end{enumerate}

\begin{enumerate}
\def\labelenumi{\alph{enumi})}
\tightlist
\item
  为使 \(u\) 是仿正规的，必须且只需
\end{enumerate}

\[
(u ^ {*}) ^ {2} \circ u ^ {2} - 2 t (u ^ {*} \circ u) + t ^ {2} 1 _ {\mathrm{E}} \geqslant 0
\]

对每个实数 t \textgreater{} 0 成立。

\begin{enumerate}
\def\labelenumi{\alph{enumi})}
\setcounter{enumi}{1}
\tightlist
\item
  为使 \(u\) 是仿正规的，必须且只需
\end{enumerate}

\[
p \circ q ^ {2} \circ p - 2 t p ^ {2} + t 1 _ {\mathrm{E}} \geqslant 0
\]

对所有 \(t > 0\) 成立，其中 \(p = (u \circ u^{*})^{1/2}\) ，\(q = (u^{*} \circ u)^{1/2}\) 。

\begin{enumerate}
\def\labelenumi{\arabic{enumi})}
\setcounter{enumi}{11}
\tightlist
\item
  \begin{enumerate}
  \def\labelenumii{\alph{enumii})}
  \tightlist
  \item
    设 u 与 v 是 E 的连续正自同态。假设对每个 t \textgreater{} 0，我们有
  \end{enumerate}
\end{enumerate}

\[
2 t u ^ {2} \circ (u ^ {2} + t ^ {2} 1 _ {\mathrm{E}}) ^ {- 1} \leqslant v \leqslant \frac {1}{2 t} (u ^ {2} + t ^ {2} 1 _ {\mathrm{E}}).
\]

证明 u = v。

\begin{enumerate}
\def\labelenumi{\alph{enumi})}
\setcounter{enumi}{1}
\tightlist
\item
  为使 \(u\) 是正规自同态，必须且只需 \(u\) 与 \(u^*\) 都是仿正规的并且有相同的核。
\end{enumerate}

\begin{enumerate}
\def\labelenumi{\arabic{enumi})}
\setcounter{enumi}{12}
\tightlist
\item
  \begin{enumerate}
  \def\labelenumii{\alph{enumii})}
  \tightlist
  \item
    设 A 是 E 的单位向量所成的无限集，它沿有限子集的余集滤子弱收敛于 0。证明存在取值于 A 的序列 \((x_{n})_{n\geqslant 1}\) 以及 E 中的标准正交序列 \((e_n)_{n\geqslant 1}\) ，使得 \(\| x_n - e_n\| \leqslant 1 / n\) 对所有 \(n\geqslant 1\) 成立 。
  \end{enumerate}
\end{enumerate}

\begin{enumerate}
\def\labelenumi{\alph{enumi})}
\setcounter{enumi}{1}
\item
  设 u 是 E 的自同态，未必连续。假设对 E 中每个标准正交序列 \((e_{n})\) ，序列 \((u(e_{n}))\) 都弱收敛于 0。证明 u 是连续且紧的。
\item
  设 u 是 E 的连续自同态。为使存在 E 的标准正交基 \((e_{n})\) 使序列 \((u(e_{n}))\) 弱收敛于 0，必须且只需：对 E 的每个自同态 v，\(v \circ u - 1_{E}\) 都不是紧的。
\end{enumerate}

\begin{enumerate}
\def\labelenumi{\arabic{enumi})}
\setcounter{enumi}{13}
\tightlist
\item
  设 X 是局部紧空间，\(\mu\) 是 X 上的正有界测度且 \(\mu(\mathrm{X}) = 1\) ，并设 \(f: X \to X\) 是 \(\mu\)-可测映射且满足 \(f(\mu) = \mu\) 。我们用 \(u_{f}\) 表示 \(\mathrm{L}_{\mathbf{C}}^{2}(\mathrm{X}, \mu)\) 的由 \(\varphi \mapsto f \circ \varphi\) 定义的酉自同态（参见 I, p. 190 的习题 16）。对每个 \(n \in N\)，我们记 \(f^{n} = f^{\circ n}\) （\(n\) 次复合），并且对 X 的每个子集 A，记 \(f^{-n}(\mathrm{A}) = g^{-1}(\mathrm{A})\) ，其中 \(g = f^{n}\) 。
\end{enumerate}

我们称 \(f\) 关于 \(\mu\) 是弱混合的，如果对 X 的所有 \(\mu\)-可测子集 A 与 B，我们有

\[
\lim _ {N \rightarrow + \infty} \frac {1}{N} \sum_ {n = 0} ^ {N} | \mu (A \cap f ^ {- n} (B) - \mu (A) \mu (B) | = 0.
\]

\begin{enumerate}
\def\labelenumi{\alph{enumi})}
\tightlist
\item
  下列条件等价：
\end{enumerate}

\begin{enumerate}
\def\labelenumi{(\roman{enumi})}
\item
  映射 f 关于 \(\mu\) 是弱混合的；
\item
  对所有 \(\varphi_{1},\varphi_{2} \in \mathrm{L}_{\mathbf{C}}^{2}(\mathrm{X}, \mu)\) ，我们有
\end{enumerate}

\[
\lim _ {N \to + \infty} \frac {1}{N} \sum_ {n = 0} ^ {N} | \langle \varphi_ {1} | u _ {f} ^ {n} (\varphi_ {2}) \rangle - \langle 1 | \varphi_ {2} \rangle \langle \varphi_ {1} | 1 \rangle | = 0;
\]

\begin{enumerate}
\def\labelenumi{(\roman{enumi})}
\setcounter{enumi}{2}
\item
  映射 \(f \times f: X \times X \to X \times X\) 是 \((\mu \otimes \mu)\) -遍历的（参见 I, p. 190 的习题 16）；
\item
  相应于特征值 1 的特征空间的维数为 1，且对每个 \(\lambda \in \mathbf{C} - \{1\}\) ，\(u_{f}\) 相应于 \(\lambda\) 的特征空间化为 0。
\end{enumerate}

（为证明 (iii) 蕴含 (i)，利用 II, p. 300 的习题 64。）

\begin{enumerate}
\def\labelenumi{\alph{enumi})}
\setcounter{enumi}{1}
\item
  考虑 I, p. 190 习题 16 中 e) 部分的例子。映射 f 关于 \(\mu\) 是弱混合的。
\item
  考虑 I, p. 190 习题 16 中问题 f) 的例子。映射 f 关于 \(\mu\) 不是弱混合的。
\end{enumerate}

\begin{enumerate}
\def\labelenumi{\arabic{enumi})}
\setcounter{enumi}{14}
\tightlist
\item
  设 \(G = (S, F, o, t, \iota)\) 是一个图（TA, II, p. 155 的定义 1）。我们假设 G 是连通的，S 是可数的，并且 S 至少有两个元素。对每个 \(x \in S\) ，我们用 \(v(x)\) 表示 G 的以 x 为端点的边的数目；我们有 \(v(x) \geqslant 1\) 对所有 \(x \in S\) 成立，并且假设 \(v(x)\) 对所有 \(x \in S\) 是有限的。对 S 中所有 x 与 y，我们用 \(a(x, y)\) 表示 G 中以 x 与 y 为端点的箭的数目。我们假设 \(a(x, y)\) 对所有 \((x, y) \in S^{2}\) 是有限的。
\end{enumerate}

我们定义 S 上的测度 \(\mu\) 如下：\(\mu(\{x\}) = v(x)\) 对所有 \(x \in S\) 成立。测度 \(\mu\) 是正的。我们用 \(\mathcal{F}(\mathrm{G})\) 表示 S 上取复值的函数所成的向量空间。我们用 \(\widetilde{\mathrm{M}}_{\mathrm{G}}\) 表示 \(\mathcal{F}(\mathrm{G})\) 的满足下式的自同态

\[
\widetilde {\mathrm{M}} _ {\mathrm{G}} (\varphi) (x) = \frac {1}{v (x)} \sum_ {y} a (x, y) \varphi (y),
\]

对所有 \(\varphi\in\mathcal{F}(\mathrm{G})\) 及每个 \(x\in\mathrm{S}\) 成立（“马尔可夫算子”）。

\begin{enumerate}
\def\labelenumi{\alph{enumi})}
\item
  自同态 \(\widetilde{\mathrm{M}}_{\mathrm{G}}\) 通过向子空间的过渡诱导出 \(\mathrm{M_G}\) ，即 \(\mathrm{L}_{\mathbf{C}}^{2}(\mathrm{S},\mu)\) 的一个自同态。
\item
  自同态 \(M_{G}\) 是自伴的。它满足
\end{enumerate}

\[
\langle \varphi | 1 - \mathrm{M} _ {\mathrm{G}} (\varphi) \rangle = \frac {1}{2} \sum_ {(x, y) \in \mathrm{S} ^ {2}} a (x, y) | \varphi (x) - \varphi (y) | ^ {2}
\]

\[
\langle \varphi | 1 + \mathrm{M} _ {\mathrm{G}} (\varphi) \rangle = \frac {1}{2} \sum_ {(x, y) \in \mathrm{S} ^ {2}} a (x, y) | \varphi (x) + \varphi (y) | ^ {2}
\]

对所有 \(\varphi\in\mathrm{L}^{2}(\mathrm{S},\mu)\) 成立。此外，我们有 \(-1\leqslant\mathrm{M}_{\mathrm{G}}\leqslant1\) 。

\begin{enumerate}
\def\labelenumi{\alph{enumi})}
\setcounter{enumi}{2}
\tightlist
\item
  设 \(d \geqslant 1\) 是整数，并假设 G 是 \(Z^{d}\) 相对于集合 \(\{e_{i}\}\) 的凯莱图，其中 \((e_{i})\) 是 \(Z^{d}\) 的典范基（TA, II, p. 221 的习题 7）。\(\varrho(\mathrm{M}_{\mathrm{G}})\) （\(M_{G}\) 的谱半径）等于 1。
\end{enumerate}

\(d\) ) 设 \(d \geqslant 3\) 是整数。假设 G 是一棵 \(d\) -正则的无限树（TA, II, p. 157），即其中每个顶点都有 \(d\) 个邻点的无限树。设 \(x_0\) 是 G 的一个固定顶点。对每个 \(\lambda \in \mathbf{C}\) ，存在唯一的函数 \(\varphi_{\lambda} \colon S \to \mathbf{C}\) ，使得 \(\varphi_{\lambda}(x_0) = 1\) 且 \(\widetilde{\mathrm{M}}_{\mathrm{G}}(\varphi_{\lambda}) = \lambda \varphi_{\lambda}\) 。此外，对每个 \(x \in S\) ，\(\varphi_{\lambda}\) 在 \(x\) 处的值只依赖于 \(x\) 到 \(x_0\) 的距离（参见 TA, II, p. 222 的习题 10）。

\begin{enumerate}
\def\labelenumi{\alph{enumi})}
\setcounter{enumi}{4}
\tightlist
\item
  \(M_{G}\) 的谱是区间 \([-2\sqrt{d-1}/d, 2\sqrt{d-1}/d]\) （“凯斯滕定理”）。特别地，我们有 \(\varrho(\mathrm{M}_{\mathrm{G}})=2\sqrt{d-1}/d\) 。
\end{enumerate}

\begin{enumerate}
\def\labelenumi{\arabic{enumi})}
\setcounter{enumi}{15}
\tightlist
\item
  我们沿用上一个习题的记号，并假设 S 与 F 是有限的。我们用 \(h_{G}\) 表示图 G 的切格常数，它定义为
\end{enumerate}

\[
h _ {\mathrm{G}} = \inf _ {\varnothing \neq \mathrm{V} \subset \mathrm{S}} \frac {\operatorname{Card} (\mathrm{E} (\mathrm{V} , \mathrm{S} - \mathrm{V}))}{\inf (\operatorname{Card} (\mathrm{V}) , \operatorname{Card} (\mathrm{S} - \mathrm{V}))},
\]

其中 E(V, W) 表示 G 的一个端点在 V 中、另一个端点在 W 中的箭所成的集。

\begin{enumerate}
\def\labelenumi{\alph{enumi})}
\tightlist
\item
  我们有 \(\varrho(\mathrm{M}_{\mathrm{G}})=1\) ，并且 \(\lambda=1\) 是 \(\mathrm{M}_{\mathrm{G}}\) 的重数为 1 的特征值。我们有 \(-1\in\mathrm{Sp}(\mathrm{M}_{\mathrm{G}})\) 当且仅当 G 是二部的（参见 TA, II, p. 219 的习题 2）；这时，\(\mathrm{M}_{\mathrm{G}}\) 的谱在 \(\lambda\mapsto-\lambda\) 下不变，并且 \(\lambda\) 的谱重数等于 \(-\lambda\) 的谱重数对所有 \(\lambda\in\mathbf{C}\) 成立。
\end{enumerate}

我们用 \(\mathrm{L}_{0}^{2}(\mathrm{S},\nu)\) 表示在 \(\mathrm{L}_{\mathbf{C}}^{2}(\mathrm{S},\nu)\) 中、\(M_{G}\) 的相应于特征值 1 与 -1 的诸特征空间之并的正交补。通过向子空间的过渡诱导的自同态 \(M_{G}\) 是 \(\mathrm{L}_{0}^{2}(\mathrm{S},\nu)\) 的自同态，记作 \(M_{G,0}\)。我们有 \(\varrho(\mathrm{M}_{\mathrm{G},0}) < 1\)。我们用 \(\lambda_{1}(\mathrm{G})\) 表示 \(1 - M_{G}\) 的最小非零特征值。

\begin{enumerate}
\def\labelenumi{\alph{enumi})}
\setcounter{enumi}{1}
\tightlist
\item
  设 \(v_{+} = \sup v(x)\) ，\(v_{-} = \inf v(x)\) 。我们有
\end{enumerate}

\[
\lambda_ {1} (\mathrm{G}) \leqslant \frac {2 v _ {+}}{v _ {-} ^ {2}} h _ {\mathrm{G}}.
\]

\begin{enumerate}
\def\labelenumi{\alph{enumi})}
\setcounter{enumi}{2}
\tightlist
\item
  设 \(\varphi\colon S\to\mathbf{R}\) 是 \(S\) 上的非常数实值函数；令 \(m=\inf\varphi(x)\) ，\(M=\sup\varphi(x)\) 。设 \(t_{0}\in\mathbf{R}\) 使得：对每个 \(t\in\mathbf{R}\) ，集合 \(W_{t}=\overline{\varphi}^{-1}(]-\infty,t_{0}[)\) 满足 \(\mathrm{Card}(W_{t})\leqslant\mathrm{Card}(S)/2\) 当且仅当 \(t<t_{0}\) 。设 \(\nu\) 是支撑含于 \([m,M]\) 的无原子测度（INT, V, p. 67 的定义 5）。令
\end{enumerate}

\[
\mathrm{A} = \frac {1}{2} \sum_ {(x, y) \in \mathrm{S} ^ {2}} a (x, y) \nu ([ \varphi (x), \varphi (y) ])
\]

\[
\mathrm{B} = \sum_ {x \in \mathrm{S}} \nu ([ t _ {0}, \varphi (x) ]),
\]

其中，当 \(a < b\) 是实数时，我们记 \(\nu([b, a]) = \nu([a, b])\) 。存在 \(t \in \mathbf{R}\) 使得

\[
\frac {\operatorname{Card} \left(\mathrm{E} \left(\mathrm{W} _ {t} , \mathrm{S} - \mathrm{W} _ {t}\right) \right.}{\inf \left(\operatorname{Card} \left(\mathrm{W} _ {t}\right) , \operatorname{Card} \left(\mathrm{S} - \mathrm{W} _ {t}\right)\right)} \leqslant \frac {\mathrm{A}}{\mathrm{B}}.
\]

\begin{enumerate}
\def\labelenumi{\alph{enumi})}
\setcounter{enumi}{3}
\tightlist
\item
  我们有
\end{enumerate}

\[
h _ {\mathrm{G}} \leqslant v _ {+} \sqrt {2 \lambda_ {1} (\mathrm{G})}.
\]

（把上一问题应用于 \(\varphi\) 取为 \(\mathrm{M_G}\) 的相应于特征值 \(\lambda_1(\mathrm{G})\) 的特征函数并且 \(\nu\) 取适当选取的测度的情形。）

\begin{enumerate}
\def\labelenumi{\alph{enumi})}
\setcounter{enumi}{4}
\tightlist
\item
  我们称一个族 \((\mathrm{G}_{j})_{j\in\mathrm{J}}\) （由有限图 \(\mathrm{G}_{j}=(\mathrm{S}_{j},\mathrm{F}_{j},o_{j},t_{j},i_{j})\) 组成）为扩张图族，如果下列条件成立：
\end{enumerate}

\begin{enumerate}
\def\labelenumi{(\roman{enumi})}
\item
  沿着 J 的有限子集的余集滤子，我们有 \(\mathrm{Card}(\mathrm{S}_{j}) \to +\infty\) ；
\item
  存在 \(C \geqslant 0\) ，使得对每个 j 及每个 \(x \in S_{j}\) ，\(G_{j}\) 在 x 处的度 \(\leqslant C\) ；
\item
  存在 \(\delta > 0\) ，使得 \(h_{\mathrm{G}_j} \geqslant \delta\) 对所有 \(j \in \mathrm{J}\) 成立。
\end{enumerate}

证明 \((\mathrm{G}_{j})\) 是扩张图族当且仅当 (i)、(ii) 以及条件

(iii') 存在 \(\delta > 0\) ，使得 \(\lambda_1(G_j) \geqslant \delta\) 对所有 \(j \in J\) 成立。

\begin{enumerate}
\def\labelenumi{\arabic{enumi})}
\setcounter{enumi}{16}
\tightlist
\item
  \begin{enumerate}
  \def\labelenumii{\alph{enumii})}
  \tightlist
  \item
    设 \(n \in \mathbf{N}\)。设 E 是维数为 \(n\) 的希尔伯特空间，\(u \in \mathcal{L}(\mathrm{E})\) ，并设 \(x \in \mathrm{E}\) 是范数为 1 的向量且满足 \(u(x) = 0\)。对每个范数为 1 的向量 \(y \in \mathrm{E}\) ，我们用 \(\Delta_u(y)\) 表示在 \(y^\circ\) 上由 \(z \mapsto \langle z | u(z) \rangle\) 定义的二次形式的判别式。
  \end{enumerate}
\end{enumerate}

对每个范数为 1 的 \(y \in E\) ，我们有 \(|\langle x | y \rangle|^{2} \Delta_{u}(x) = \Delta_{u}(y)\) 。（注意，\(\Delta_{u}(y) = \langle a | \bigwedge u^{n-1}(a) \rangle / \langle a | a \rangle\) 对空间 \(\bigwedge^{n-1} y^{\circ} \subset \bigwedge^{n-1} E\) 的每个非零元素 a 成立，其中空间 \(\bigwedge^{n-1} E\) 赋有其典范希尔伯特结构，参见 EVT, V, p.33；然后分别在 x 与 y 正交以及 x = y 的情形验证该公式。）

\begin{enumerate}
\def\labelenumi{\alph{enumi})}
\setcounter{enumi}{1}
\tightlist
\item
  设 A 是复系数的 \((n,n)\) 型矩阵，满足 \({}^{t}\overline{A}=A\) 。对每个满足 \(1\leqslant i\leqslant n\) 的整数 i，用 \(A_{i}\) 表示 \((n-1,n-1)\) 型的、通过删除 A 的第 \(i\) 行与第 \(i\) 列所得的矩阵。设 \((e_{1},\ldots,e_{n})\) 是 \(C^{n}\) 的、赋有其典范希尔伯特结构且由 A 的特征向量组成的标准正交基，并用 \(\lambda_{i}\) 表示 u 相应于 \(e_{i}\) 的特征值。记 \(e_{i}=(e_{i,j})_{1\leqslant j\leqslant n}\) 。
\end{enumerate}

对所有 \((i,j)\) ，我们有

\[
\left| e _ {i, j} \right| ^ {2} \prod_ {k \neq i} \left(\lambda_ {i} - \lambda_ {k}\right) = \det \left(\lambda_ {i} - A _ {j}\right).
\]

（化归到 \(\lambda_{i}=0\) 的情形，然后把上一问题应用于空间 \(C^{n}\) 、自同态 \(u\colon y\mapsto Ay\) 以及向量 \(x=e_{i}\) 。）

\begin{enumerate}
\def\labelenumi{\arabic{enumi})}
\setcounter{enumi}{17}
\tightlist
\item
  设 X 是紧可度量化空间，\(\mu\) 是 X 上的正测度。记 \(\mathrm{E} = \mathrm{L}_{\mathbf{C}}^{2}(\mathrm{X}, \mu)\)。设 \((\mathcal{X}_{n})_{n \in \mathbf{N}}\) 是把 X 分成可测集的有限分划的序列，使得分划 \(X_{n+1}\) 比 \(X_{n}\) 更细对所有 \(n \in N\) 成立，并且使得 \(X_{n}\) 的各部分的直径中的最大者当 \(n \to +\infty\) 时趋于 0 。
\end{enumerate}

设 \(E_{n} \subset E\) 是由满足如下条件的 \(f \in \mathrm{L}^{2}(\mathrm{X}, \mu)\) 组成的空间：f \(\mu\)-几乎处处在 \(X_{n}\) 的每个部分 Y 上为常数。我们令 \(E_{-1} = \{0\}\)。对每个 \(n \geqslant -1\)，用 \(p_{n}\) 表示 E 的以 \(E_{n}\) 为像的正交投影算子。

\begin{enumerate}
\def\labelenumi{\alph{enumi})}
\item
  序列 \((p_{n})_{n\in\mathbf{N}}\) 在赋有逐点收敛拓扑的 E 的自同态所成的空间 \(\mathcal{L}_{s}(\mathrm{E})\) 中收敛于恒等映射。（首先考虑 \(p_{n}(f)\) 当 \(f\in\mathcal{C}(\mathrm{X})\) 时的极限。）
\item
  对每个 \(n \in N\) ，自同态 \(q_{n} = p_{n} - p_{n-1}\) 是到 \(F_{n} = E_{n} \cap E_{n-1}^{\circ}\) 上的正交投影算子，并且诸空间 \(F_{n}\) 两两正交。
\item
  我们有 \(\sum_{n} q_{n} = 1_{\mathrm{E}}\) 在 \(\mathcal{L}_s(\mathrm{E})\) 中成立，并且空间 E 是诸空间 \(\mathrm{F}_n\) 的希尔伯特直和。
\end{enumerate}

\begin{enumerate}
\def\labelenumi{\arabic{enumi})}
\setcounter{enumi}{18}
\tightlist
\item
  设 \(\mathrm{X} \subset \mathbf{R}\) 是 \(\mathbf{R}\) 的紧子集，\(\mu\) 是 \(\mathrm{X}\) 上的正测度，而 \(u\) 是由 \(\mathrm{X}\) 的恒等函数给出的乘法自同态。
\end{enumerate}

我们记 \(\mathrm{E}=\mathrm{L}^{2}(\mathrm{X},\mu)\) ，并用 \(\mathcal{L}_{s}(\mathrm{E})\) 表示赋有逐点收敛拓扑的 E 的自同态所成的空间。

设 \(\varepsilon > 0\)。取 \(r > 0\) 使得 \(\mathrm{X} \subset ] - r, r[\) ，并取 \(\ell \geqslant 1\) 使得 \(r2^{-\ell} < \varepsilon\) 。对每个整数 \(n \in \mathbf{N}\) 以及每个整数 \(j\) ，其中 \(|j| \leqslant 2^{n + \ell}\) ，令

\[
\mathrm{X} _ {n, j} = \mathrm{X} \cap \Big ] \frac {r j}{2 ^ {n + \ell}}, \frac {r (j + 1)}{2 ^ {n + \ell}} \Big ]
\]

\begin{enumerate}
\def\labelenumi{\alph{enumi})}
\item
  对每个整数 \(n \in N\)，族 \(\mathcal{X}_{n} = (\mathrm{X}_{n,j})_{j}\) 是 X 的有限分划；诸集合 \(X_{n,j}\) 是可测的且直径 \(\leqslant r2^{-n-\ell}\) 。对每个 \(n \in N\)，分划 \(X_{n+1}\) 比分划 \(X_{n}\) 更细。因此我们可以对 \((\mathcal{X}_{n})\) 应用上一个习题的结果并沿用其记号，特别是空间 \(E_{n}\) 与 \(F_{n}\) 以及正交投影算子 \(p_{n}\) 与 \(q_{n}\) 。诸空间 \(E_{n}\) 是有限维的。
\item
  级数
\end{enumerate}

\[
v = \sum_ {n \in \mathbf {N}} q _ {n} u q _ {n}
\]

在空间 \(\mathcal{L}_{s}(\mathrm{E})\) 中收敛；存在 E 的标准正交基 B 使得 \(v \in \mathcal{D}_{\mathrm{B}}(\mathrm{E})\) 。

\begin{enumerate}
\def\labelenumi{\alph{enumi})}
\setcounter{enumi}{2}
\tightlist
\item
  自同态 \(w = v - u\) 是紧的。（考虑乘法自同态 \(u_n\) ，它由函数 \(g_n \in \mathrm{E}_n\) 给出，使得对每个整数 \(j\) （\(|j| \leqslant 2^{n + \ell}\) ）以及所有 \(x \in \mathrm{X}_{n,j}\) ，有 \(g_n(x) = rj2^{-n - \ell}\) ；然后验证 \(\| u - u_n \| \leqslant r2^{-n - \ell}\) 以及 \(u_n p_n = p_n u_n\) 。）
\end{enumerate}

\begin{enumerate}
\def\labelenumi{\arabic{enumi})}
\setcounter{enumi}{19}
\tightlist
\item
  设 E 是可数型复希尔伯特空间，\(u \in \mathcal{L}(\mathrm{E})\) 是埃尔米特自同态。设 \(\varepsilon\) 是严格正的实数。存在 E 的标准正交基 B、埃尔米特对角自同态 \(v \in \mathcal{D}_{\mathrm{B}}(\mathrm{E})\) 以及紧自同态 \(w \in \mathcal{L}^{\mathrm{c}}(\mathrm{E})\) ，使得 \(u = v + w\) 且 \(\|w\| < \varepsilon\) （首先考虑 u 有循环向量的情形，然后应用上一个习题。）
\end{enumerate}

21) 设 E 是可数型复希尔伯特空间，\(u \in \mathcal{L}(\mathrm{E})\) 是正规自同态。设 \(\varepsilon\) 是严格正的实数。存在 E 的标准正交基 B、对角自同态 \(v \in \mathcal{D}_{\mathrm{B}}(\mathrm{E})\) 以及紧自同态 \(w \in \mathcal{L}^{\mathrm{c}}(\mathrm{E})\) ，使得 \(u = v + w\) 且 \(\|w\| < \varepsilon\) 。

\begin{enumerate}
\def\labelenumi{\arabic{enumi})}
\setcounter{enumi}{21}
\tightlist
\item
  设 G 是局部紧交换拓扑群，\(\mu\) 是 G 上的哈尔测度。
\end{enumerate}

设 \(u \in \mathcal{L}(\mathrm{L}^2(\mathrm{G},\mu))\) 满足：\(u\) 与乘法算子 \(m_{\chi}\) 对每个特征标 \(\chi \in \widehat{\mathrm{G}}\) 交换。证明存在 \(g \in \mathrm{L}^{\infty}(\mathrm{G},\mu)\) 使得 \(u = m_g\) 。

\begin{enumerate}
\def\labelenumi{\arabic{enumi})}
\setcounter{enumi}{22}
\tightlist
\item
  设 \(u\) 是 E 的自同态，并设 A 是 \(\mathcal{L}(\mathrm{E})\) 的由 \(u\) 生成的闭含单位元子代数。
\end{enumerate}

若 v 是 E 的自同态且 \(u \circ v - v \circ u\) 属于 A，则 \(u \circ v - v \circ u\) 是拟幂零的。特别地，若 E 非零，我们有 \(u \circ v - v \circ u \neq 1_{E}\) 。（利用 I, p. 167 的习题 10；注意该结果对任何巴拿赫空间 E 都成立。）

\begin{enumerate}
\def\labelenumi{\arabic{enumi})}
\setcounter{enumi}{23}
\tightlist
\item
  设 \(\mathrm{A} \subset \mathcal{L}(\mathrm{E})\) 是 \(\mathcal{L}(\mathrm{E})\) 的含单位元的自伴子代数。我们用 \(\mathrm{A}''\) 表示 A 的双交换子。我们用 \(\mathcal{L}(\mathrm{E})_s\) 表示赋有逐点收敛拓扑的空间 \(\mathcal{L}(\mathrm{E})\) 。
\end{enumerate}

\begin{enumerate}
\def\labelenumi{\alph{enumi})}
\item
  空间 \(A''\) 是 \(\mathcal{L}(\mathrm{E})\) 的含单位元的自伴子代数，并且它在 \(\mathcal{L}(\mathrm{E})_s\) 中是闭的。
\item
  设 \(v \in A''\) 且 \(x \in E\)。向量 \(v(x)\) 属于诸 \(u(x)\) （\(u \in A\)）所成之集在 E 中的闭包。（设 F 是该闭包；注意 E 的以 F 为像的正交投影算子属于 \(A'\)。）
\item
  代数 \(A''\) 与 A 在 \(\mathcal{L}(\mathrm{E})_s\) 中的闭包重合。（对 E 的元素的每个有限族 \((x_i)_{i \in \mathrm{I}}\) ，把上一问题应用于空间 \(\mathrm{E}^{\mathrm{I}}\) 、像 \(\widetilde{\mathrm{A}}\) （即 A 在含单位元对合态射 \(f: \mathcal{L}(\mathrm{E}) \to \mathcal{L}(\mathrm{E}^{\mathrm{I}})\) （其中 \(f(u)(y_i) = (u(y_i))_{i \in \mathrm{I}}\) ）下的像），以及元素 \((x_i)\) （它属于 \(\mathrm{E}^{\mathrm{I}}\) ）。）
\end{enumerate}

\(A''\) 与 A 在 \(\mathcal{L}(\mathrm{E})_s\) 中闭包的相等就是冯·诺伊曼双交换子定理。

\begin{enumerate}
\def\labelenumi{\arabic{enumi})}
\setcounter{enumi}{24}
\tightlist
\item
  设 U 是 R 的开子集。普遍可测函数 \(f: U \to R\) 称为勒夫纳函数，如果对每个复希尔伯特空间 E 以及 E 的满足 \(\operatorname{Sp}(u) \subset \operatorname{U}\) 与 \(\operatorname{Sp}(v) \subset \operatorname{U}\) 的埃尔米特自同态 u 与 v，条件 \(u \leqslant v\) 在 \(\mathcal{L}(\operatorname{E})\) 中蕴含 \(f(u) \leqslant f(v)\) 。我们用 \(\operatorname{Loe}(\operatorname{U})\) 表示 U 上勒夫纳函数所成的集。
\end{enumerate}

\begin{enumerate}
\def\labelenumi{\alph{enumi})}
\item
  设 \(f \in \text{Loe}(U)\)。函数 f 在 U 上递增。常值函数 1 与 U 的恒等函数都属于 \(\text{Loe}(U)\) 。
\item
  设 \(f \in \text{Loe}(\mathbf{R}_+^*)\)。设 E 是复希尔伯特空间，u 是 E 的正自同态。对 E 的每个正交投影算子 p 以及每个 \(\varepsilon > 0\) ，我们有 \(f(u)_p \leqslant f((1 + \varepsilon)u_p)\) ，其中用 \(v_p\) 表示 \(\text{Im}(p)\) 的由 \(x \mapsto p(v(x))\) 定义的自同态（对所有 \(v \in \mathcal{L}(\text{E})\) ）。（把 E 等同于希尔伯特直和 \(\text{Im}(p) \oplus \text{Im}(1_{\text{E}} - p)\) ，并将 u 表示为矩阵 \(\begin{pmatrix} a & b \\ c & d \end{pmatrix}\)。注意 \(u \leqslant \begin{pmatrix} (1 + \varepsilon)a & 0 \\ 0 & (1 + \varepsilon^{-1})d \end{pmatrix}\) 对所有 \(\varepsilon > 0\) 成立。）
\item
  设 \(f \in \text{Loe}(\mathbf{R}_{+}^{*})\)。函数 f 在 \(R_{+}^{*}\) 上是凹的（设 v 是 \(C^{2}\) 的在典范基下为对角的自同态，并设 \(g \in \text{SO}(\mathbf{R}^{2})\) 视为 \(C^{2}\) 的自同态。把上一问题应用于以 \(Ce_{1}\) 为像的正交投影算子且取 \(u = g^{*}vg\) 。）
\item
  集合 \(\operatorname{Loe}(\mathbf{R})\) 是形如 \(f(x)=ax+b\) （其中 \((a,b)\in\mathbf{R}_{+}\times\mathbf{R}\) ）的函数所成的集。（把上一问题既应用于 f，也应用于函数 \(x\mapsto-f(-x)\) 。）
\item
  设 \(t \in \mathbf{R}_{+}^{*}\) 。函数 \(f: x \mapsto -(t + x)^{-1}\) 属于 \(\operatorname{Loe}(\mathbf{R}_{+}^{*})\)。
\end{enumerate}

\(f)\) 对任何有界正测度 \(\nu\) （在 \(\mathbf{R}_{+}^{*}\) 上且使函数 \(x\mapsto x^{-1}\) 为 \(\nu\) -可积），函数 \(f\colon\mathbf{R}_{+}^{*}\to\mathbf{R}\) 由

\[
f (x) = \int_ {\mathbf {R} _ {+} ^ {*}} \frac {x}{x + t} d \nu (t)
\]

所给出的函数属于 \(\text{Loe}(\mathbf{R}_{+}^{*})\)；它在 \(R_{+}^{*}\) 上可微，并且还满足下列条件：

\begin{enumerate}
\def\labelenumi{(\roman{enumi})}
\item
  \(\lim_{t\to0}f(t)=0;\)
\item
  \(\lim_{t\to0}f'(t)\) 在 R 中存在；
\item
  f 在 \(+\infty\) 处有有限极限。
\end{enumerate}

\begin{enumerate}
\def\labelenumi{\arabic{enumi})}
\setcounter{enumi}{25}
\tightlist
\item
  设 E 是复希尔伯特空间，a 是 E 的正自同态。我们记 \(\|x\|_{a} = \langle x | a(x) \rangle\) ；它是 E 上的半范数。我们同样用 \(\|u\|_{a}\) 表示 E 赋有该半范数时其自同态 u 的半范数。
\end{enumerate}

我们称函数 \(f \colon \mathbf{R}_+^* \to \mathbf{R}_+^*\) 为插值函数，如果对每个复希尔伯特空间 E、E 的每个正自同态 \(a\) 以及每个 \(u \in \mathcal{L}(\mathrm{E})\) ，条件 \(\| u \| \leqslant 1\) 与 \(\| u \|_a \leqslant 1\) 蕴含 \(\| f(u) \|_a \leqslant 1\) 。

\begin{enumerate}
\def\labelenumi{\alph{enumi})}
\item
  函数 \(f \colon \mathbf{R}_+^* \to \mathbf{R}_+^*\) 是插值函数，当且仅当对每个复希尔伯特空间 E、E 的每个正自同态 \(a\) 以及每个 \(u \in \mathcal{L}(\mathrm{E})\) ，条件 \(u^*u \leqslant 1_{\mathrm{E}}\) 与 \(u^*au \leqslant a\) 蕴含 \(u^*f(a)u \leqslant f(a)\) 。
\item
  \(R_{+}^{*}\) 上的任何勒夫纳函数都是插值函数。（给定 E 以及 \(u \in \mathcal{L}(E)\) （范数 \(\leqslant 1\) ），令 \(w = \sqrt{1_{E} - u^{*}u}\) ，\(w' = \sqrt{1_{E}-uu^{*}}\) ；考虑 \(F = E \oplus E\) ，并注意由 \(v(x,y) = (u(x)+w'(y), -w(x)+u(y))\) 定义的 F 的自同态 v 是等距的；然后考虑 \(a' \colon (x,y) \mapsto (a(x),0)\) ，并把上一个习题的 b) 应用于以 \(E \oplus \{0\}\) 为像的正交投影算子。）
\item
  每个插值函数都是 \(\mathbf{R}_{+}^{*}\) 上的勒夫纳函数。（设 \(0 \leqslant u \leqslant v\) 是复希尔伯特空间 E 的自同态且 \(v\) 可逆；令 \(\mathrm{F} = \mathrm{E} \oplus \mathrm{E}\) ，考虑 F 的自同态 \(a = u \oplus v\) 以及自同态 \(w\) ，它由矩阵 \(\left( \begin{array}{cc}0 & 1_{\mathrm{E}}\\ v^{-1/2}u^{1/2} & 0 \end{array} \right)\) 表示；证明 \(w^{*}w \leqslant 1_{\mathrm{F}}\) 且 \(w^{*}aw \leqslant a\) 。）
\item
  设 f 是插值函数。设 E 是复希尔伯特空间；我们记 \(\mathcal{R}(u)=\frac{1}{2}(u+u^{*})\) 对所有 \(u\in\mathcal{L}(\mathrm{E})\) 。对 \(\mathcal{L}(\mathrm{E})\) 中的 u 与 v，证明：\(e^{-tv^{*}}\circ u\circ e^{-tv}\leqslant u\) 对所有 \(t\in R_{+}\) 成立当且仅当 \(\mathcal{R}(u\circ v)\geqslant0\) 。（对给定的 \(x\in E\)，计算函数 \(t\mapsto\langle e^{-tv}x|(ue^{-tv})x\rangle\) 的导数。）
\end{enumerate}

由此推出：对 E 的所有自同态 \(u\) 与 \(v\) ，当条件

\[
u \geqslant 0, \quad \mathcal {R} (v) \geqslant 0, \quad \mathcal {R} (u v) \geqslant 0
\]

得到满足时，我们有 \(\mathcal{R}(f(u)v) \geqslant 0\) 。

27) 我们打算证明习题 25 问题 \(b\) 中论断的逆命题。为此，将利用下面 V, p. 503 习题 23 的结果与记号。特别地，G 表示群 \(\mathbf{R}_{+}^{*}\) ，它赋有哈尔测度 \(\mu = x^{-1}dx\) ，并与群 \(i\mathbf{R}\) 通过映射 \((x, it) \mapsto x^{it}\) 对偶。我们用 X 表示 G 的恒等映射。

我们分别用 co、si 与 ta 表示 iR 上的函数 \(it \mapsto \cos(i\pi t)\)、\(it \mapsto \sin(i\pi t)\) 与 \(it \mapsto \tan(i\pi t)\) 。

设 \(f \in \text{Loe}(\mathbf{G})\) 满足习题 25 问题 \(f\) 的条件 (i)、(ii) 与 (iii)。

\begin{enumerate}
\def\labelenumi{\alph{enumi})}
\item
  设 \(\alpha\) 是由 \(x \mapsto x / (1 + x)\) 给出的 \(G\) 上的函数。我们有 \(\mathrm{I}_{\alpha} = ] - 1,0[\) ，并且 \(\widehat{\alpha}(s) = -\pi / \sin(\pi s)\) 对 \(s \in \mathrm{B}_f\) 成立。
\item
  函数 ta 在 \(i\mathbf{R}\) 上有界。我们令 \(v = F \circ m_{ta} \circ \overline{F}\) ；它是 \(L^{2}(G)\) 的自同态，满足 \(\|v\| \leqslant 1\) 且 \(\mathcal{R}(v) = 0\) 。
\item
  设 \(\varepsilon > 0\) 是实数。我们用 \(u_{\varepsilon}\) 表示由 \(h_{\varepsilon} = \mathrm{X}/(1 + \varepsilon\mathrm{X})\) 给出的 \(\mathrm{L}^2(\mathrm{G})\) 上的乘法自同态。我们有 \(u_{\varepsilon} \geqslant 0\) ，并且对每个 \(\varphi \in \mathrm{L}^2(\mathrm{G})\) 以及每个有界连续函数 \(f\) （在 \(\mathrm{G}\) 上），有
\end{enumerate}

\[
\langle \varphi | \mathcal {R} (f (u _ {\varepsilon}) v) \varphi \rangle = \frac {1}{2} \langle f \circ h _ {\varepsilon} | \mathrm{Q} (\varphi) \rangle
\]

\(\mathrm{Q}(\varphi) = \overline{\varphi} v(\varphi) + \overline{v(\varphi)}\varphi .\)

\begin{enumerate}
\def\labelenumi{\alph{enumi})}
\setcounter{enumi}{3}
\item
  设 F 是 \(\mathrm{L}^{2}(i\mathbf{R})\) 的由诸函数 \(t \mapsto e^{\alpha(it - s_{0})^{2}}\) （其中 \(\alpha \in R_{+}^{*}\) ，\(s_{0} \in C\) ）生成的向量子空间；它在 \(\mathrm{L}^{2}(i\mathbf{R})\) 中稠密。
\item
  设 \(\psi \in \mathrm{F}\) 。存在 \(\varphi \in \mathrm{L}^2 (\mathbf{G})\) 使得 \(\widehat{\varphi} = \mathrm{co}\cdot \psi\)。令 \(g = \mathrm{Q}(\varphi)\)。我们有 \(g\in \mathrm{L}^{1}(\mathrm{G})\) ，并且存在 \(\mathbf{C}\) 上的全纯函数，它到 \(i\mathbf{R}\) 的限制与 \(\widehat{g}\) 一致。此外
\end{enumerate}

\[
\widehat {g} = \frac {1}{2 i \pi} (\psi^ {*} * \psi) \mathrm{si}.
\]

\(f)\) 对每个 \(\varepsilon > 0\) ，我们有 \(\mathcal{R}(u_{\varepsilon} \circ v) \geqslant 0\) 。（计算 \(\langle \varphi | \mathcal{R}(u_{\varepsilon} u) \varphi \rangle\) ，当 \(\varphi\) 是上一问题中那样的函数时；并像 V, p. 503 习题 23 的问题 \(d\) 中那样使用 G 上的卷积。）

\(g)\) 设 \(f \in \mathrm{Loe}(\mathbf{R}_{+}^{*})\)。设 \(g\) 是与 \(\psi \in F\) 如问题 \(e)\) 中那样相伴的函数。我们则有

\[
\int_ {\mathrm{G}} f g d \mu \geqslant 0
\]

以及

\[
\int_ {\mathbf {R}} \widehat {f} (- \sigma - i t) \sin (\pi (\sigma + i t)) (\psi^ {*} * \psi) (\sigma + i t) d t \geqslant 0
\]

对所有 \(\sigma\in]0,1[\) 成立。

\begin{enumerate}
\def\labelenumi{\alph{enumi})}
\setcounter{enumi}{7}
\tightlist
\item
  设 \(\sigma \in ]0,1[\) 。函数 \(q\) 由 \(q(s) = \widehat{f}(-s)\sin(\pi s)\) （对 \(s\) 满足 \(\mathcal{R}(s) = \sigma\) ）定义，它满足
\end{enumerate}

\[
\sum_ {i = 0} ^ {n} \sum_ {j = 0} ^ {n} \bar {t} _ {i} t _ {j} q \left(\frac {1}{2} (\bar {z} _ {i} + z _ {j})\right) \geqslant 0
\]

对所有 \(n \in N\) 以及所有复数族 \((t_{i})_{0 \leqslant i \leqslant n}\) 与 \((z_{i})_{0 \leqslant i \leqslant n}\) （它们满足对每个 i 有 \(\mathcal{R}(z_{i}) = \sigma\) ）成立。

\begin{enumerate}
\def\labelenumi{\roman{enumi})}
\tightlist
\item
  存在 G 上的正测度 \(\nu_0\) ，使得 \(\widehat{\nu}_0(it) = \widehat{f}(-it)\sin(i\pi t)/\pi\) 对所有 \(t\in\mathbf{R}\) 成立；我们有 \(\mathrm{X}^{\sigma}\in\mathrm{L}^{1}(\mathrm{G},\nu_{0})\) ，若 \(0<\sigma<1\) 。
\end{enumerate}

\begin{enumerate}
\def\labelenumi{\alph{enumi})}
\setcounter{enumi}{9}
\tightlist
\item
  设 \(\nu = \check{\nu}_0\)。测度 \(\nu\) 是有界的。（先验证 \(]-1,0[ \subset I_{\nu}\) ，再验证 \(\widehat{\nu}(-t)\) 当 \(t \to 0\) 时有界，以此证明 \(\nu([1,+\infty[)\) 有限。）对每个 \(x\) ，我们有
\end{enumerate}

\[
f (x) = \int_ {\mathrm{G}} \frac {x}{x + t} d \nu (t).
\]

\begin{enumerate}
\def\labelenumi{\arabic{enumi})}
\setcounter{enumi}{27}
\tightlist
\item
  \begin{enumerate}
  \def\labelenumii{\alph{enumii})}
  \tightlist
  \item
    设 \(D \subset C\) 是开单位圆盘。设 f 是 \(\overline{D}\) 的一个邻域中的全纯函数的芽。对每个 \(z \in D\) ，我们有
  \end{enumerate}
\end{enumerate}

\[
f (z) = i \mathcal {I} (f (0)) + \frac {1}{2 \pi} \int_ {0} ^ {2 \pi} \mathcal {R} (f (e ^ {i \theta})) k (e ^ {i \theta}, z) d \theta
\]

其中 \(k(z, w) = (z + w) / (z - w)\) 对 \((z, w) \in \mathrm{D} \times \mathrm{D}\) 。（注意

\[
k (e ^ {i \theta}, z) = 1 + 2 \sum_ {n = 1} ^ {+ \infty} z ^ {n} e ^ {- i n \theta}
\]

并使用 f 在 0 附近的泰勒展开。）

\begin{enumerate}
\def\labelenumi{\alph{enumi})}
\setcounter{enumi}{1}
\tightlist
\item
  设 \(f\) 是 D 上的全纯函数，使得 \(\mathcal{R}(f(z)) > 0\) 对所有 \(z \in \mathrm{D}\) 成立。测度 \(\mu_r = \frac{1}{2\pi} \mathcal{R}(f(re^{i\theta})) d\theta\) （\([0, 2\pi]\) 上的）当 \(r \to 1\) 时淡收敛于一个有界正测度 \(\nu\) ，并且我们有
\end{enumerate}

\[
f (z) = i \mathcal {I} (f (0)) + \frac {1}{2 \pi} \int_ {0} ^ {2 \pi} \mathcal {R} (f (e ^ {i \theta})) k (e ^ {i \theta}, z) d \nu (\theta)
\]

对所有 \(z \in D\) 成立。（化归到 \(f(0) = 1\) 的情形，并注意这时诸测度 \(\mu_{r}\) 是总质量为 1 的正测度。）

\begin{enumerate}
\def\labelenumi{\alph{enumi})}
\setcounter{enumi}{2}
\tightlist
\item
  设 H 是由这样的 \(z \in C\) 组成的 C 的开子集：\(\mathcal{I}(z) > 0\)。设 \(f: H \to C\) 是全纯函数。f 的像含于 H 当且仅当存在 R 上的有界正测度 \(\nu\) 以及实数 \(a \geqslant 0\) 使得
\end{enumerate}

\[
f (z) = \mathcal {R} (f (i)) + a z + \int_ {\mathbf {R}} \frac {1 + x z}{x - z} d \nu (x)
\]

对所有 \(z \in \mathbf{H}\) 成立。

\begin{enumerate}
\def\labelenumi{\alph{enumi})}
\setcounter{enumi}{3}
\tightlist
\item
  设 I 是 \(\mathbf{R}\) 中的非空开区间，并设 \(f\colon \mathrm{I}\to \mathbf{R}\) 。记 \(\mathrm{J} = \mathbf{R} - \mathrm{I}\) 。下列条件等价（“勒夫纳定理”）：
\end{enumerate}

\begin{enumerate}
\def\labelenumi{(\roman{enumi})}
\item
  函数 \(f\) 属于 \(\operatorname{Loe}(\mathbf{I})\) ；
\item
  存在解析函数 \(\widetilde{f}\) ，它定义在 \((\mathbf{C} - \mathbf{R}) \cup I\) 上、到 I 的限制与 \(f\) 一致，并且满足 \(\mathcal{I}(f(z)) > 0\) 若 \(\mathcal{I}(z) > 0\) ；
\item
  存在 J 上的有界测度 \(\nu\) 以及偶 \((a,b)\in\mathbf{R}_{+}\times\mathbf{R}\) 使得
\end{enumerate}

\[
f (x) = a x + b + \int_ {\mathrm{J}} \frac {1 + x y}{y - x} d \nu (y)
\]

对所有 \(x \in I\) 成立。

（首先直接验证 (iii) 蕴含 (i) 与 (ii)；为证明 (i) 蕴含 (iii)，单独处理 I = R 的情形，并把其他情形化归为 I = R \(_{+}^{*}\) 的情形，对此应用上一个习题；最后，为证明 (ii) 蕴含 (iii)，应用 c) 部分并验证由此给出的 R 上的测度 \(\nu\) 的支撑含于 J。）

